\section{Discussion}
\label{sec:discussion}

\subsection{Principal findings}

Four statements are supported by the evidence, listed in order of how well it supports them.

\begin{enumerate}
  \item \textbf{Muscle loss during bed rest follows a curved, not linear, path}: about 3\pp{} per doubling of days
        (95\% CI 2.5 to 3.4), with the reference curve at $-12\%$ by day 60 and $-15\%$ by day 119. The corpus
        cannot say \emph{which} curve: logarithmic, saturating and spline forms are indistinguishable.
  \item \textbf{Which muscle is measured matters more than anything else in the dataset.} Thirty per cent of the
        variance lies between muscles within a campaign; plantar flexors lose 5.7\pp{} more than knee extensors
        and 5.6\pp{} more than their antagonists, the dorsiflexors; antigravity extensors lose roughly twice what
        flexors lose. The flexible models of tier~2 lean on the same two variables.
  \item \textbf{At 32 campaigns, tabular machine learning adds nothing to a simple curve.} The four families find
        the muscle structure---pooled $R^2$ doubles or triples---but do not convert it into lower out-of-campaign
        error. That is a statement about how much signal the published literature contains, not about the
        algorithms.
  \item \textbf{A model that reads the other campaigns' data in context adds a little.} The language-model forecast
        is 0.42\pp{} more accurate than the curve, survives the removal of study-identifying detail and every
        presentation check, and is not explained by recognition of published campaigns. It was not pre-registered,
        it misses its declared bar and its ranges are overconfident; it is a finding to report, not the headline.
\end{enumerate}

\subsection{The duration--response}

The clearest single result is negative about shape and positive about magnitude. A straight line in days is the worst
of the three baseline forms, and the curvature is real: roughly 40\% of the loss the reference curve reaches by day
119 has already occurred by day 14 ($-6.0\%$ of $-15.3\%$). What the data cannot do is choose between a curve that
keeps falling with the logarithm of time and one that approaches a plateau. The 32 campaigns are concentrated at 5--21
and 56--90 days, with a single campaign at 119 days, and both forms pass through that support equally well. The
saturating form was kept as the headline by a rule declared in advance, because its parameters are interpretable---a
time constant of 60 days and an eventual loss of about 17\%---not because the data prefer it. For planning purposes
the honest summary is the curve and its band over the observed range, not the asymptote.

Two properties of the curve should be kept in view when quoting it. It is drawn for one scenario: a control arm, knee
extensors, measured as CT cross-sectional area. For MRI volume it sits about half a percentage point higher, and for
plantar flexors about 5.7\pp{} lower. And it is fitted to group means: it describes what a group of similar people
would lose on average, not what an individual will lose.

\subsection{Muscle-specific vulnerability: mechanical role}

The ranking reproduces, as adjusted coefficients with intervals, the selective vulnerability reported by
muscle-resolved bed-rest studies~\citep{belavy2009,miokovic2012,marusic2021}: the plantar flexors are the most
affected family and the hip rotators and adductors the least. The corpus also offers a partial test of the two usual
explanations. Tibialis anterior, the principal dorsiflexor, is slow-fibre dominant like the soleus~\citep{johnson1973}
but does not carry body weight; if fibre type governed atrophy, the dorsiflexors should lose like the plantar flexors.
They do not: they lose like the knee extensors ($-10.2\%$ against $-10.1\%$ at day 60) and 5.6\pp{} less than the
plantar flexors. The comparison rests on eight dorsiflexor rows from four campaigns and on a classification that has
not yet been signed off, so it is suggestive rather than decisive, but it points towards mechanical role.

The countermeasure effect, 3.7\pp{} less loss in countermeasure arms, pools every countermeasure type, from flywheel
exercise to nutrition. It shows that the dataset can detect a countermeasure effect at all; it cannot say which
countermeasure works, which would need more campaigns per modality than the literature provides.

\subsection{Why flexible learners did not help}

The tier-2 null result is the most important methodological finding of the project, and the explanation is the
sample size. There are 346 rows, but 32 independent campaigns, and the error that matters is the error on a campaign
the model has never seen. Most of what distinguishes one campaign from another---its imaging protocol, its
population, the exact definition of its muscles, its countermeasure dose---is either not recorded in a way that can be
modelled or is perfectly confounded with the campaign itself. The flexible families learn the muscle structure, which
raises pooled $R^2$, but the part of the residual that decides out-of-campaign error is between-campaign, and no
column explains it. With roughly a tenth of the sample size at which flexible learners begin to behave
well~\citep{vanderploeg2014}, this is the expected outcome, and it was prepared for in advance.

The answer this gives to the question ``why not simply use machine learning?'' is a better one than any accuracy
figure would have been: because at this sample size the unit of evidence is the campaign, and the literature holds 32
of them. It also explains why the meta-regression, and not a predictive model, is the right primary tool: it
estimates exactly the structure the learners found, with intervals, and it does not need to predict unseen campaigns
to be useful.

\subsection{What tier 3 does and does not show}

The language-model forecast is the only method in the project whose advantage over the curve has a paired interval
that excludes zero. The ablations establish what the gain does not need---anything that identifies a study---and that
it depends on the reference values: shuffle them and the gain becomes a loss. The recognition probe finds that the
model can name some campaigns, but not the ones carrying the gain. The likeliest reading is that the model does a kind
of case-based reasoning that tier 2 could not: it picks out the reference rows that match the target---same muscle,
same kind of group, a nearby day---and weighs them, where tier 2 received the same information as one-hot columns and
a duration basis. That is an interpretation, not a tested mechanism.

The limits are as clear as the gain. The method was added after the null result; it misses the 15\% bar declared for
it; its ranges are overconfident, so its uncertainty cannot be quoted; its answers vary slightly between identical
requests; and on the history arm it does not use the campaign's earlier scans at all---the benefit of history comes
entirely from the anchor the code supplies. Pretraining exposure to some of these papers cannot be excluded by any
cross-validation.

\subsection{Implications for mission planning}

The results support three statements a countermeasure planner can use, at group level and within the observed range
of 5 to 119 days: that most of the loss seen by four months has accumulated within the first weeks; that the calf's
antigravity muscles are the ones to protect first; and that countermeasure arms in the literature lose measurably
less than controls. They do not support a prediction for an individual astronaut, a statement about actual
microgravity, or a number for day 180. Extrapolation to a six-month transit was specified but has not been run; when
it is, it will be computed from the tier-1 curve only, reported as a prediction interval drawn on the same axes as the
data, and accompanied by the statement that the longest observation is 119 days and that the remainder is an
assumption about curve shape.

\subsection{Reconciliation with the accepted abstract}
\label{sec:abstract-reconciliation}

The accepted abstract was written before the literature search and is the contract for the presentation. The results
confirm its central claims and correct several of its numbers (\cref{tab:abstract}). None of the corrections weakens
the claims; most strengthen the evidence behind them.

\begin{table}[htbp]
\centering
\caption[Reconciliation with the accepted abstract]{The accepted abstract against the results.}
\label{tab:abstract}
\small
\begin{tabularx}{\linewidth}{@{}L{4.6cm} Y@{}}
\toprule
Abstract & Result \\
\midrule
``15 bed-rest and HDBR studies (1992--2025)'' & 52 studies (1992--2026) from 36 independent campaigns; the modelling
  subset holds 41 studies from 32 campaigns \\
``unloading durations of 14--119 days'' & Planned durations and scan days both span 5--119 days \\
Variables include participant characteristics & Recorded, but not modelled: age and sex are almost perfectly confounded
  with the campaign and are reserved for sensitivity analyses S8 and S9 \\
Linear regression, random forest, SVR, gradient boosting & As stated; linear regression is penalised (ridge) \\
Leave-one-study-out cross-validation & Leave-one-\emph{campaign}-out: stricter, because papers from one campaign share
  participants \\
SHAP-based interpretation & Done for the best family; the top two features are stable, the third is not \\
``consistent relationship between duration and loss'' & Confirmed: curved, about $-3$\pp{} per doubling of days \\
Soleus, gastrocnemius and triceps surae most susceptible, ``losses approaching 20--30\% during prolonged bed rest'' &
  Plantar flexors are the most affected family, $-15.8\%$ at day 60 by the model; individual group means reach $-30\%$
  at the longest durations, but the modelled averages do not \\
``Exercise countermeasures consistently associated with reduced muscle loss'' & Countermeasure arms lose 3.7\pp{} less
  (95\% CI 1.5 to 5.9), all countermeasure types pooled \\
Duration and muscle-specific vulnerability as major determinants & Confirmed by tier 1 (variance components,
  coefficients) and by tier-2 importance \\
\bottomrule
\end{tabularx}
\end{table}

\section{Limitations}
\label{sec:limitations}

The limitations are stated in full; several are properties of the literature rather than of the analysis.

\begin{enumerate}
  \item \textbf{Small effective sample.} The modelling subset holds 32 independent campaigns. Every interval is wide
        accordingly, and the covariate budget is about three study-level terms.
  \item \textbf{Clustering and dominance.} One campaign (MEDES LTBR) supplies 40\% of the full dataset and a quarter of
        the modelling subset; 15 campaigns contribute five rows or fewer. Campaign weighting and campaign-level
        validation mitigate this, but the must-show sensitivity analysis that drops the largest campaign (S2) has not
        been run.
  \item \textbf{Search coverage.} The search starts in 2013; older work enters only through a convenience corpus of
        nine papers. Embase was not searched, 1,023 records were not screened in the time available, and 19 included
        studies report their results only as charts without a baseline.
  \item \textbf{No human double extraction.} The extraction was verified against the sources with AI assistance, and
        the pre-search corpus rows were not verified by a second person at all.
  \item \textbf{Heterogeneous measurements in one target.} MRI volume, CT and MRI cross-sectional area, DXA lean mass
        and ultrasound thickness are treated as one target with a modality term; the MRI-only sensitivity analysis
        (S4) is outstanding. The measurement site along a muscle, which matters~\citep{miokovic2012}, is ignored in
        the primary model, and 547 rows do not state which leg was imaged.
  \item \textbf{Mixed granularity and two estimators.} Composites and components coexist, resolved by a rule whose
        opposite (S1) has not been run; the target mixes the printed mean of individual changes with a value
        recomputed from group means.
  \item \textbf{Publication bias} is plausible in this literature and cannot be quantified with 32 campaigns.
  \item \textbf{Bed rest is an analogue.} Nothing here establishes that the coefficients transfer to microgravity;
        the 14 spaceflight rows are too few to test it and eight of them are lumbar multifidus.
  \item \textbf{No data beyond 119 days}, and a single campaign at 119 days.
  \item \textbf{Mostly young men.} 347 of the 425 lower-limb unloading rows come from men-only groups and 38 from
        older adults, in six campaigns.
  \item \textbf{Group means, not individuals.} Every row is an arm mean; nothing can be said about individual
        susceptibility.
  \item \textbf{The NASA campaign is thin.} Three of its four in-bed-rest occasions rest on one participant.
  \item \textbf{Unsigned classification.} The muscle classification awaits physiological sign-off, a blocking
        condition for quoting coefficients by class.
  \item \textbf{Assumptions of the model.} Campaigns are assumed independent of one another; missing dispersions and
        ages are assumed uninformative; the tier-2 families were fitted unweighted.
  \item \textbf{Tier 3} is post hoc, overconfident, not bit-reproducible, and exposed to possible pretraining on the
        source papers.
  \item \textbf{Most sensitivity analyses are outstanding}, including the three declared as must-show.
\end{enumerate}

\section{Status of the work and remaining tasks}
\label{sec:outstanding}

The internal schedule placed the modelling hard stop on 21 September 2026; the report and slides follow, with the
upload deadline on 10 October and the presentation on 15 or 16 October. The following items remain, in order of
priority.

\begin{enumerate}
  \item \textbf{Sensitivity analyses S1, S2 and S4} (composite-first resolution, without MEDES LTBR, MRI volume only),
        declared as must-show in this report; then the remaining analyses where they would change a conclusion.
  \item \textbf{Physiological sign-off} of \file{data/muscle\_map.csv} and of the fitted relationships by the
        scientific lead---a blocking criterion, and a line edit in one file.
  \item \textbf{The 180-day extrapolation} with a prediction interval, as a clearly caveated backup result.
  \item \textbf{Recalibration of tier 3's ranges} by a rule declared in advance, if its uncertainty is to be quoted.
  \item \textbf{Reconcile the PRISMA counts} at the full-text stage with the final extraction.
  \item \textbf{Close P3 and P4 formally}: tag the results and merge the branch stack ending at \ReportBranch{} into
        \code{main}.
  \item \textbf{Repository housekeeping before publication}: add a licence (code and dataset need different ones),
        remove the full-text digests that reproduce copyrighted text, acknowledge NASA as a data source, check the
        terms under which tier 3's cached answers may be published, and obtain both leads' agreement to the one
        schema change in v1.1.
  \item \textbf{Two data checks}: two countermeasure-arm rows carry \code{cm\_modality = none}, and the figure-derived
        extraction table fails the validator unless run in partial mode (23 rows lack sex or age).
  \item \textbf{Author instructions} from the organisers (slide format, time limit, disclosure requirements), still
        outstanding.
\end{enumerate}

\section{Conclusions}
\label{sec:conclusions}

The published bed-rest literature, organised into one frozen dataset of 742 comparisons from 52 studies and 36
campaigns, supports a quantified duration--response---about 3 percentage points of muscle per doubling of days of
unloading, curved rather than linear---and a muscle-specific vulnerability led by the plantar flexors, with
countermeasure arms losing measurably less. Tested honestly, by holding out whole campaigns, four machine-learning
families add nothing to a simple duration curve, because the sample size is the number of campaigns, not the number of
rows; a language model reading the other campaigns' data adds a small, robust but post-hoc improvement. The dataset
and the framework are the contribution, and every number in this report can be regenerated from them with one command.
